\chapter{Introduction}

\section{Background}
Confidentiality is the first pillar of the CIA triad and the historical starting
point of cryptography. Classical ciphers such as Caesar and Playfair illustrate
substitution and digraph techniques, while Simplified DES (S-DES) is a scaled-down
teaching version of the Feistel block-cipher structure used by DES and, in spirit,
by AES. Studying all three inside a single working network application makes the
transition from paper exercises to real message and file protection concrete: the
same plaintext must be encoded, framed, transmitted over a socket, received,
decrypted and verified byte-for-byte.

\section{Problem Statement}
A plain TCP socket transfers data in cleartext, so any party on the network path
can read messages and files. The problem addressed by this prototype is therefore:
\emph{how can a client and a server exchange text messages and bulk binary files
over TCP such that only the endpoints holding the shared key can read the
content, while the mechanism itself remains simple enough to inspect and reason
about?}

\section{Objectives}
\begin{enumerate}
  \item Implement Caesar, Playfair and byte-wise S-DES encryption and decryption
        as a reusable Python module.
  \item Build a TCP server that accepts an encrypted frame, decrypts it, and
        displays the ciphertext alongside the recovered plaintext.
  \item Build an interactive client that lets the user choose the algorithm and
        key, encrypt either a typed message or a 1~MB binary file, and decrypt the
        server's reply.
  \item Verify correctness by comparing the recovered file with the original
        byte-for-byte, and measure the throughput of the S-DES implementation.
  \item Analyse the residual security weaknesses of the prototype and recommend
        defensible alternatives.
\end{enumerate}

\section{Scope and Assumptions}
The prototype is an educational artefact. The shared key is agreed out of band
(the user types it into both endpoints) and is carried inside the frame header for
convenience; no key exchange, certificate validation, integrity check or replay
protection is implemented. S-DES operates one byte at a time in Electronic
Codebook (ECB) mode. Caesar and Playfair operate on text only, and Playfair
additionally discards non-alphabetic characters. Testing was carried out on the
loopback interface with one client connection at a time, because the server
handles connections sequentially.

\section{Development Environment}
\begin{table}[!ht]
\centering
\caption{Prototype implementation and test environment}
\label{tab:env}
\begin{tabular}{|l|l|}
\hline
\textbf{Item} & \textbf{Value} \\ \hline
Language & Python 3.14.3 (standard library only) \\ \hline
Modules used & \texttt{socket}, \texttt{json}, \texttt{os}, \texttt{math} \\ \hline
Transport & TCP, stream sockets (\texttt{AF\_INET}, \texttt{SOCK\_STREAM}) \\ \hline
Server bind address & \texttt{0.0.0.0}, port \texttt{9000} \\ \hline
Client target & \texttt{127.0.0.1}, port \texttt{9000} \\ \hline
Ciphers & Caesar, Playfair, S-DES (byte-wise, ECB) \\ \hline
Test data & \texttt{sample\_1mb.bin} (1{,}048{,}576 B), \texttt{sample\_10kb.bin} (10{,}240 B) \\ \hline
Source repository & \url{https://github.com/Trinadh-31/IAS} \\ \hline
\end{tabular}
\end{table}

\chapter{System Design}

\section{Architecture Overview}
The prototype follows a three-layer split: a cipher library, two network
endpoints, and a sample-data generator. Both endpoints import the same cipher
functions, so encryption and decryption are guaranteed to be symmetric.
Figure~\ref{fig:arch} shows the arrangement and the direction of data flow.

\begin{figure}[!ht]
    \centering
    \begin{tikzpicture}[
        font=\small,
        node distance=6mm,
        mod/.style={draw=blue!55!black, fill=blue!5, rounded corners=2pt,
                    align=center, inner sep=5pt, minimum height=20mm, text width=32mm},
        lib/.style={mod, draw=orange!65!black, fill=orange!8, minimum height=16mm, text width=33mm},
        dat/.style={mod, draw=green!45!black, fill=green!6, minimum height=12mm, text width=34mm},
        arr/.style={-{Latex[length=2mm]}, blue!55!black, thick},
        darr/.style={-{Latex[length=2mm]}, gray, dashed}
    ]
    \node[mod] (cli) at (0,0) {\textbf{client.py}\\[2pt]\scriptsize Interactive CLI\\ \scriptsize selects algorithm + key\\ \scriptsize encrypts and frames data};
    \node[mod] (srv) at (8.4,0) {\textbf{server.py}\\[2pt]\scriptsize TCP listener :9000\\ \scriptsize parses frame, decrypts\\ \scriptsize prints and replies};
    \node[lib] (lib) at (4.2,2.3) {\textbf{ciphers.py}\\[2pt]\scriptsize caesar / playfair / sdes\\ \scriptsize encrypt and decrypt};
    \node[dat] (dat) at (4.2,-2.4) {\textbf{Sample data}\\[2pt]\scriptsize sample\_1mb.bin\\ \scriptsize sample\_10kb.bin};

    \draw[arr] (cli.north) -- node[above left=-1mm and -2mm]{\scriptsize import} (lib.west);
    \draw[arr] (srv.north) -- node[above right=-1mm and -2mm]{\scriptsize import} (lib.east);
    \draw[arr] ([yshift=3mm]cli.east) -- node[above]{\scriptsize encrypted request} ([yshift=3mm]srv.west);
    \draw[arr] ([yshift=-4mm]srv.west) -- node[below]{\scriptsize encrypted reply} ([yshift=-4mm]cli.east);
    \draw[darr] (dat.west) -- (cli.south);
    \draw[darr] (dat.east) -- (srv.south);
    \end{tikzpicture}
    \caption{Component architecture of the client--server cipher prototype; both endpoints share one cipher module and communicate over a TCP socket on port 9000.}
    \label{fig:arch}
\end{figure}

\section{Module Responsibilities}
Table~\ref{tab:modules} lists each source file and the responsibility it owns.
Keeping the ciphers in one module avoids duplicated key-schedule logic and makes
the algorithms independently testable.

\begin{table}[!ht]
\centering
\caption{Source files and their responsibilities}
\label{tab:modules}
\begin{tabular}{|l|l|}
\hline
\textbf{File} & \textbf{Responsibility} \\ \hline
\texttt{ciphers.py} & Caesar, Playfair and S-DES primitives; self-test in \texttt{\_\_main\_\_} \\ \hline
\texttt{server.py} & Listener, frame parsing, decryption, display, encrypted reply \\ \hline
\texttt{client.py} & Menu-driven UI, key entry, encryption, framing, reply handling \\ \hline
\texttt{make\_files.py} & Generates the 1~MB and 10~KB random sample files \\ \hline
\end{tabular}
\end{table}

\section{Message Frame Format}
Because TCP is a byte stream and not a message-oriented protocol, the prototype
defines its own application-layer framing, shown in Figure~\ref{fig:frame}. A
four-byte big-endian length prefix tells the peer how many bytes of JSON header
follow; the header in turn declares \texttt{payload\_size}, so the receiver knows
exactly how many ciphertext bytes to read. The helper \texttt{recv\_all()} loops
until the requested number of bytes has arrived, which prevents the short-read
bug that naive \texttt{recv()} calls suffer on large transfers.

\begin{figure}[!ht]
    \centering
    \begin{tikzpicture}[font=\small,
        fld/.style={draw=blue!55!black, align=center, minimum height=14mm, inner sep=4pt}]
    \node[fld, fill=blue!8, text width=20mm] (a) at (0,0) {\textbf{Header}\\ \textbf{length}\\[2pt]\scriptsize 4 bytes\\ \scriptsize big-endian};
    \node[fld, fill=orange!10, text width=45mm, right=0.6mm of a] (b) {\textbf{JSON header}\\[2pt]\scriptsize mode, algorithm, key,\\ \scriptsize filename, payload\_size};
    \node[fld, fill=green!8, text width=38mm, right=0.6mm of b] (c) {\textbf{Payload}\\[2pt]\scriptsize ciphertext bytes\\ \scriptsize (message or file)};
    \draw[{Latex[length=2mm]}-{Latex[length=2mm]}, gray] ([yshift=-3mm]a.south west) -- ([yshift=-3mm]c.south east);
    \node[gray, below=6mm of b.south, text width=0.95\linewidth, align=center]
        {\scriptsize read in order by \texttt{recv\_all()}: 4 bytes, then \texttt{hdr\_len} bytes, then \texttt{payload\_size} bytes};
    \end{tikzpicture}
    \caption{Application-layer frame: four-byte length prefix, JSON header and ciphertext payload.}
    \label{fig:frame}
\end{figure}

\begin{table}[!ht]
\centering
\caption{Fields carried in the JSON header}
\label{tab:header}
\begin{tabular}{|l|l|l|}
\hline
\textbf{Field} & \textbf{Example} & \textbf{Purpose} \\ \hline
\texttt{mode} & \texttt{message} / \texttt{file} & Selects the receiver's handling path \\ \hline
\texttt{algorithm} & \texttt{caesar} / \texttt{playfair} / \texttt{sdes} & Selects the decryption routine \\ \hline
\texttt{key} & \texttt{3} / \texttt{SECURE} / \texttt{1010000010} & Shared key for that algorithm \\ \hline
\texttt{filename} & \texttt{sample\_1mb.bin} & Output name for file mode \\ \hline
\texttt{payload\_size} & \texttt{1048576} & Exact ciphertext length in bytes \\ \hline
\end{tabular}
\end{table}

\section{Control Flow}
The client presents a menu with three actions: send a message, send the 1~MB file
using S-DES, or quit. For a message the user picks one of the three algorithms and
supplies the matching key; the client prints the ciphertext, opens a fresh
connection, transmits the frame and waits for the reply. The server runs an
accept loop: for every connection it reads one frame, decrypts it, prints both
representations, and then responds. For Caesar and Playfair the response is an
encrypted acknowledgement string; for S-DES the response is an encrypted 10~KB
file, which exercises the bulk path in the opposite direction.

\chapter{Cryptographic Algorithms}

\section{Caesar Cipher}
The Caesar cipher shifts every alphabetic character by a fixed integer, wrapping
within its case and leaving digits, spaces and punctuation untouched. Decryption
reuses the encryption routine with the additive inverse of the shift,
$(-k) \bmod 26$, which halves the code that has to be trusted.

\begin{algorithm}[H]
\SetAlgoLined
\KwIn{plaintext $P$, shift $k$}
\KwOut{ciphertext $C$}
$C \leftarrow \varepsilon$\;
\ForEach{character $c \in P$}{
  \uIf{$c$ is lowercase}{
    append $\mathrm{chr}((\mathrm{ord}(c)-\mathrm{ord}(\text{`a'})+k) \bmod 26 + \mathrm{ord}(\text{`a'}))$ to $C$\;
  }
  \uElseIf{$c$ is uppercase}{
    append $\mathrm{chr}((\mathrm{ord}(c)-\mathrm{ord}(\text{`A'})+k) \bmod 26 + \mathrm{ord}(\text{`A'}))$ to $C$\;
  }
  \Else{
    append $c$ unchanged to $C$\;
  }
}
\Return $C$\;
\caption{Caesar encryption as implemented in \texttt{encrypt\_caesar()}}
\end{algorithm}

\section{Playfair Cipher}
Playfair encrypts digraphs using a $5 \times 5$ key matrix. The implementation
normalises the text to uppercase letters, maps \texttt{J} to \texttt{I}, inserts
\texttt{X} between repeated letters of a pair, and pads an odd-length string with a
trailing \texttt{X}. Three substitution rules then apply to each pair: same row
shifts right, same column shifts down, otherwise the letters are replaced by the
opposite corners of the rectangle they define.

\begin{algorithm}[H]
\SetAlgoLined
\KwIn{plaintext $P$, keyword $K$}
\KwOut{ciphertext $C$}
$T \leftarrow$ normalise and digraph-pad $P$ (uppercase, $J \rightarrow I$, insert $X$)\;
$M \leftarrow$ $5\times5$ matrix from unique letters of $K$, then remaining alphabet without $J$\;
\For{$i \leftarrow 0$ \KwTo $|T|-1$ \textbf{step} $2$}{
  $(r_a,c_a) \leftarrow \mathrm{pos}(T_i)$; \quad $(r_b,c_b) \leftarrow \mathrm{pos}(T_{i+1})$\;
  \uIf{$r_a = r_b$}{
    append $M[r_a][(c_a{+}1) \bmod 5]$ and $M[r_b][(c_b{+}1) \bmod 5]$\;
  }
  \uElseIf{$c_a = c_b$}{
    append $M[(r_a{+}1) \bmod 5][c_a]$ and $M[(r_b{+}1) \bmod 5][c_b]$\;
  }
  \Else{
    append $M[r_a][c_b]$ and $M[r_b][c_a]$\;
  }
}
\Return $C$\;
\caption{Playfair encryption as implemented in \texttt{encrypt\_playfair()}}
\end{algorithm}

\section{Simplified DES}
S-DES is a two-round Feistel cipher with an 8-bit block and a 10-bit key. The key
schedule permutes the key with P10, splits it into two halves, and applies
circular left shifts of one and two positions before compressing each result with
P8 to produce subkeys $K_1$ and $K_2$. Each round applies the function $f_K$:
expand/permute the right half with EP, XOR with the subkey, substitute through
S-boxes $S_0$ and $S_1$, permute with P4, and XOR into the left half. The block is
swapped between rounds, and the inverse initial permutation ends the process.
Decryption is identical with the subkeys applied in reverse order.

\begin{table}[!ht]
\centering
\caption{Permutation tables used by the S-DES implementation}
\label{tab:sdes}
\begin{tabular}{|l|l|}
\hline
\textbf{Table} & \textbf{Values} \\ \hline
P10 & 3, 5, 2, 7, 4, 10, 1, 9, 8, 6 \\ \hline
P8 & 6, 3, 7, 4, 8, 5, 10, 9 \\ \hline
IP & 2, 6, 3, 1, 4, 8, 5, 7 \\ \hline
IP$^{-1}$ & 4, 1, 3, 5, 7, 2, 8, 6 \\ \hline
EP & 4, 1, 2, 3, 2, 3, 4, 1 \\ \hline
P4 & 2, 4, 3, 1 \\ \hline
\end{tabular}
\end{table}

\begin{algorithm}[H]
\SetAlgoLined
\KwIn{plaintext byte $b$, 10-bit key $K$}
\KwOut{ciphertext byte $b'$}
$(K_1,K_2) \leftarrow$ \texttt{generate\_subkeys}$(K)$ \tcp*{P10, shifts, P8}
$x \leftarrow \mathrm{IP}(\mathrm{bits}(b))$\;
$x \leftarrow f_K(x, K_1)$ \tcp*{EP, XOR, $S_0/S_1$, P4}
$x \leftarrow \mathrm{swap}(x)$ \tcp*{exchange 4-bit halves}
$x \leftarrow f_K(x, K_2)$\;
$b' \leftarrow \mathrm{byte}(\mathrm{IP}^{-1}(x))$\;
\Return $b'$\;
\caption{S-DES single-block encryption (\texttt{sdes\_encrypt\_block()})}
\end{algorithm}

\section{Key Material}
Table~\ref{tab:keys} summarises the key formats and the size of the resulting
search space, which determines how easily each cipher can be brute-forced.

\begin{table}[!ht]
\centering
\caption{Key formats and effective key space per algorithm}
\label{tab:keys}
\begin{tabular}{|l|l|p{6.0cm}|}
\hline
\textbf{Algorithm} & \textbf{Key format} & \textbf{Effective key space} \\ \hline
Caesar & integer shift, e.g.\ \texttt{3} & 25 useful keys \\ \hline
Playfair & keyword, e.g.\ \texttt{SECURE} & $25! \approx 1.55\times10^{25}$ matrices, far fewer in practice \\ \hline
S-DES & 10-bit string, e.g.\ \texttt{1010000010} & $2^{10} = 1{,}024$ keys \\ \hline
\end{tabular}
\end{table}

\chapter{Implementation}

\section{Cipher Module}
The cipher module exposes symmetric pairs of functions and a bytes-level wrapper
for S-DES that iterates over the input one byte at a time. The listing below shows
the key schedule and the bulk wrapper.

\begin{singlespace}
\begin{lstlisting}[language=Python, caption=S-DES key schedule and bulk encryption (\texttt{ciphers.py})]
def _generate_sdes_subkeys(key10):
    p10 = [3,5,2,7,4,10,1,9,8,6]
    p8  = [6,3,7,4,8,5,10,9]
    k = _permute(key10, p10)
    left, right = k[:5], k[5:]
    left, right = _left_shift(left,1), _left_shift(right,1)
    k1 = _permute(left+right, p8)
    left, right = _left_shift(left,2), _left_shift(right,2)
    k2 = _permute(left+right, p8)
    return k1, k2

def sdes_encrypt_bytes(data: bytes, key10: str) -> bytes:
    if len(key10) != 10:
        raise ValueError('SDES key must be 10 bits')
    key_bits = [int(c) for c in key10]
    out = bytearray()
    for b in data:
        out.append(sdes_encrypt_block(b, key_bits))
    return bytes(out)
\end{lstlisting}
\end{singlespace}

\section{Client}
The client builds the frame and writes it in three ordered sends, then reads the
reply with the same discipline.

\begin{singlespace}
\begin{lstlisting}[language=Python, caption=Framing and transmission on the client (\texttt{client.py})]
def send_header_and_payload(conn, header: dict, payload: bytes):
    hdr_b = json.dumps(header).encode()
    conn.send(len(hdr_b).to_bytes(4,'big'))
    conn.send(hdr_b)
    if payload:
        conn.send(payload)

header = {'mode':'message','algorithm':'sdes','key':key,
          'payload_size':len(payload)}
\end{lstlisting}
\end{singlespace}

\section{Server}
On the server side \texttt{recv\_all()} guarantees complete reads, and the
handler dispatches on \texttt{mode} and \texttt{algorithm}.

\begin{singlespace}
\begin{lstlisting}[language=Python, caption=Complete reads and frame parsing on the server (\texttt{server.py})]
def recv_all(conn, n):
    buf = b''
    while len(buf) < n:
        chunk = conn.recv(n - len(buf))
        if not chunk:
            raise ConnectionError('socket closed')
        buf += chunk
    return buf

hdr_len = int.from_bytes(recv_all(conn, 4), 'big')
hdr = json.loads(recv_all(conn, hdr_len).decode())
payload = recv_all(conn, hdr.get('payload_size', 0))
\end{lstlisting}
\end{singlespace}

\section{Bulk File Handling}
\begin{sloppypar}
In file mode the client reads \texttt{sample\_1mb.bin}, encrypts the whole buffer
with S-DES and sends it as a single payload. The server stores the ciphertext as
\texttt{received\_sample\_1mb.bin}, writes the decrypted result as
\texttt{decrypted\_sample\_1mb.bin}, and then encrypts and returns
\texttt{sample\_10kb.bin}. The client decrypts that response into
\texttt{server\_sent\_sample\_10kb.bin}. This produces four artefacts that can be
compared with the originals to prove that no bytes were lost or altered.
\end{sloppypar}

\section{Execution Procedure}
\begin{singlespace}
\begin{lstlisting}[language=bash, caption=Running the prototype]
python make_files.py     # generate sample_1mb.bin and sample_10kb.bin
python server.py         # terminal 1: listener on port 9000
python client.py         # terminal 2: interactive menu
# SDES key must be exactly 10 bits, e.g. 1010000010
\end{lstlisting}
\end{singlespace}

\chapter{Testing and Results}

\section{Methodology}
Each cipher was exercised in message mode, and S-DES was additionally exercised in
file mode with the 1~MB sample. Correctness was judged on two criteria: the
plaintext printed by the receiver had to match what was sent, and the decrypted
file on disk had to be byte-identical to the source file. Timings were taken
around the bulk encryption and decryption calls on the same machine that hosted
both endpoints.

\section{Functional Test Cases}
Table~\ref{tab:tests} lists the executed cases with the values actually observed
in the terminals.

\begin{table}[!ht]
\centering
\caption{Functional test cases and observed results}
\label{tab:tests}
\resizebox{\textwidth}{!}{%
\begin{tabular}{|l|l|l|l|l|}
\hline
\textbf{\#} & \textbf{Case} & \textbf{Key} & \textbf{Observed output} & \textbf{Result} \\ \hline
1 & Caesar message & 3 & \texttt{Lqirupdwlrq Dvvxudqfh dqg Vhfxulwb} & Pass \\ \hline
2 & Caesar reply decrypt & 3 & \texttt{Server reply: received your message} & Pass \\ \hline
3 & Playfair message & \texttt{SECURE} & \texttt{ISKYIQEWUMKC} & Pass \\ \hline
4 & Playfair reply decrypt & \texttt{SECURE} & \texttt{SERVERREPLYRECEIVEDYOURMESSAGE} & Pass \\ \hline
5 & S-DES message \texttt{HELLO} & \texttt{1010000010} & ciphertext \texttt{e050a5a547}, server printed \texttt{b'HELLO'} & Pass \\ \hline
6 & S-DES 1~MB file transfer & \texttt{1010000010} & \texttt{decrypted\_sample\_1mb.bin} identical, 1{,}048{,}576 B & Pass \\ \hline
7 & S-DES 10~KB reply file & \texttt{1010000010} & \texttt{server\_sent\_sample\_10kb.bin} identical, 10{,}240 B & Pass \\ \hline
8 & Playfair round trip of \texttt{HELLO WORLD} & \texttt{SECURE} & recovered \texttt{HELXLOWORLDX} & Partial \\ \hline
9 & Invalid S-DES key length & \texttt{101} & \texttt{ValueError: SDES key must be 10 bits} & Pass \\ \hline
\end{tabular}%
}
\end{table}

\section{Console Transcripts}
The following transcript was captured on the server while the Caesar and S-DES
message tests ran.

\begin{singlespace}
\begin{lstlisting}[caption=Server output during message tests]
Starting server on 0.0.0.0 9000
Client connected ('127.0.0.1', 55270)
Received mode=message algorithm=caesar key=3 size=34
Ciphertext (client->server): Lqirupdwlrq Dvvxudqfh dqg Vhfxulwb
Decrypted plaintext: Information Assurance and Security
Sent reply ciphertext and plaintext shown above (server side)
Client connected ('127.0.0.1', 46710)
Received mode=message algorithm=sdes key=1010000010 size=5
Received SDES ciphertext (hex, first 200 bytes): e050a5a547
Decrypted plaintext (first 200 bytes): b'HELLO'
Server will send SDES-encrypted 10KB file back to client
Sent encrypted 10KB file to client
\end{lstlisting}
\end{singlespace}

\begin{singlespace}
\begin{lstlisting}[caption=Client output during the 1~MB S-DES file transfer]
Enter 10-bit SDES key to use: 1010000010
Encrypted 1MB file preview (hex, first 200 bytes): d7207a0cefabd6c927ac9d5b...
Received file header from server: {'mode': 'file', 'algorithm': 'sdes',
 'key': '1010000010', 'filename': 'sample_10kb.bin', 'payload_size': 10240}
Saved decrypted file from server to server_sent_sample_10kb.bin
\end{lstlisting}
\end{singlespace}

\section{Performance}
Table~\ref{tab:perf} reports the measured cost of the byte-wise S-DES
implementation. The throughput of roughly 0.13~MB/s is a direct consequence of
processing eight bits per Python function call; a table-driven or block-based
implementation, or a call into a compiled library, would raise this by orders of
magnitude. Caesar and Playfair were applied only to short messages, where the cost
was not measurable at this resolution.

\begin{table}[!ht]
\centering
\caption{Measured S-DES performance on the test machine}
\label{tab:perf}
\begin{tabular}{|l|l|l|l|}
\hline
\textbf{Payload} & \textbf{Encrypt} & \textbf{Decrypt} & \textbf{Throughput} \\ \hline
1~MB (1{,}048{,}576 B) & 7.58~s & 7.57~s & $\approx$ 0.13~MB/s \\ \hline
10~KB (10{,}240 B) & 0.073~s & 0.073~s & $\approx$ 0.13~MB/s \\ \hline
5~B message (\texttt{HELLO}) & $<$ 0.001~s & $<$ 0.001~s & not measurable \\ \hline
\end{tabular}
\end{table}

\section{Observations and Known Defects}
\begin{enumerate}
  \item \textbf{Playfair is not loss-free.} Because the preparation step strips
        spaces and inserts filler letters, \texttt{HELLO WORLD} decrypts to
        \texttt{HELXLOWORLDX}. This is inherent to Playfair, but it means the
        cipher cannot be used for binary or exact-text payloads.
  \item \textbf{Asymmetric reply in S-DES message mode.} When a message is sent
        with S-DES the server answers with a \texttt{file} frame, while the client
        only prints \texttt{Server sent a file back; use receive logic}. The reply
        path for that combination is unfinished.
  \item \textbf{Key echoed in the header.} Every frame carries the key in
        cleartext, so an observer on the path recovers the plaintext trivially.
  \item \textbf{Single-threaded server.} Connections are handled one at a time, so
        a long 1~MB transfer blocks all other clients.
  \item \textbf{No padding or ECB alternative.} S-DES encrypts each byte
        independently, so repeated bytes are visible in the ciphertext.
\end{enumerate}

